\chapter{Устойчивость и погрешность операторной модели}
\label{rig:chapter}
\label{op:chapter}

Пусть $p=(A,n,k,\theta)$ есть набор параметров. При единственности
решения усечённой краевой задачи ему соответствует рассеянное поле
$u_s(p)$. Так определяется оператор решения $\mathcal S:p\mapsto u_s(p)$.
Он вообще нелинеен по параметрам среды, хотя при фиксированных $A,n,k$
задача линейна по нагрузке. Конечно-элементный метод задаёт дискретный
оператор $\mathcal S_h:p\mapsto U(p)$, а обученная модель приближает
его значения. Возникают два вопроса. Насколько велика ошибка $U(p)$
относительно $u_s(p)$? Насколько близок к $U(p)$ прогноз по сокращённому
набору коэффициентов? Для ответов нужны разные нормы и разные константы.

\section{Устойчивость дискретной краевой задачи}
\label{rig:spaces}
\label{rig:discrete}

Пусть $\Omega$ есть ограниченная липшицева область, $k>0$, а $A,n$ ---
вещественные измеримые функции, для которых почти всюду
\begin{equation}
 0<A_{\min}\le A\le A_{\max}<\infty,\qquad
 0<n_{\min}\le n\le n_{\max}<\infty.\label{rig:coeff-bounds}
\end{equation}
Как и в главе~\ref{math:fem}, примем $A=1$ в окрестности
$\Gamma=\partial\Omega$. Введём пространство тестовых функций и норму
\begin{equation}
 V=H^1(\Omega;\C),\qquad \|v\|_{V,k}^2=\|\nabla v\|_{L^2}^2+k^2\|v\|_{L^2}^2.
 \label{rig:vnorm}
\end{equation}
В пространстве $V'$ непрерывных сопряжённо-линейных функционалов
норма нагрузки равна $\|\ell\|_{V',k}=\sup_{v\ne0}|\ell(v)|/\|v\|_{V,k}$.
Слабая форма и определяемый ею оператор имеют вид
\begin{equation}
 (\mathcal A_k w)(v)=a_k(w,v)
 =\int_\Omega A\nabla w\cdot\nabla\bar v
 -k^2\int_\Omega nw\bar v-\ii k\int_\Gamma w\bar v.\label{rig:form}
\end{equation}
Для решения выполняется $\mathcal A_k u_s=\ell$. Из неравенств
Коши--Буняковского и следа $\|v\|_{L^2(\Gamma)}\le C_{\rm tr}\|v\|_{H^1}$
следует оценка
\begin{equation}
 |a_k(w,v)|\le C_a(k)\|w\|_{V,k}\|v\|_{V,k},\qquad
 C_a(k)=A_{\max}+n_{\max}+k C_{\rm tr}^2\max(1,k^{-2}).\label{rig:continuity}
\end{equation}
С другой стороны, для ненулевой постоянной функции
$\operatorname{Re}a_k(v,v)=-k^2\int_\Omega n|v|^2<0$. Стандартное
доказательство устойчивости через положительную вещественную часть
здесь неприменимо.

Пусть $V_h\subset V$ есть пространство непрерывных кусочно линейных
функций на согласованной невырожденной триангуляции. В этой главе $u_{s,h}$
обозначает решение при точных интегралах. Тогда дискретная задача
задаётся равенством
\begin{equation}
 a_k(u_{s,h},v_h)=\ell(v_h)\quad(v_h\in V_h).\label{rig:galerkin}
\end{equation}
Нижнюю границу величины действия формы на единичном векторе задаёт число
\begin{equation}
 \beta_h(k)=\inf_{w_h\ne0}\sup_{v_h\ne0}
 \frac{|a_k(w_h,v_h)|}{\|w_h\|_{V,k}\|v_h\|_{V,k}},\qquad w_h,v_h\in V_h.
 \label{rig:beta}
\end{equation}
Константа зависит также от $A,n$. Положительность $\beta_h$ называется
дискретным условием inf-sup. В квадратной конечномерной системе она
равносильна обратимости.

\begin{theorem}[Оценка решения и ошибки по невязке]
При $\beta_h>0$ решение существует и единственно для каждой нагрузки.
Для приближения $v\in V_h$ положим $r_h(w_h)=\ell(w_h)-a_k(v,w_h)$.
Тогда
\begin{align}
 \|u_{s,h}\|_{V,k}&\le\beta_h^{-1}\|\ell\|_{V_h',k},\label{rig:discrete-stability}\\
 \|u_{s,h}-v\|_{V,k}&\le\beta_h^{-1}\|r_h\|_{V_h',k}.\label{rig:residual-estimate}
\end{align}
\end{theorem}
\begin{proof}
По \eqref{rig:beta} нулевой нагрузке может отвечать только нулевое
решение. Поэтому у квадратной системы тривиальное ядро и она обратима.
Подстановка решения в определение $\beta_h$ даёт первое неравенство.
Для второго достаточно применить его к задаче с решением $u_{s,h}-v$
и нагрузкой $r_h$.
\end{proof}

Пусть $L_h$ есть матрица системы \eqref{math:fem-system}, $K_0$ ---
матрица жёсткости при $A=1$ и $H_h=K_0+k^2M$. Вектор узловых значений
$z$ задаёт $\|z_h\|_{V,k}^2=z^*H_hz$. Для вектора нагрузки $F$
двойственная норма равна $(F^*H_h^{-1}F)^{1/2}$. Переход к переменным
$H_h^{1/2}z$ в \eqref{rig:beta} даёт
\begin{equation}
 \beta_h(k)=\sigma_{\min}(H_h^{-1/2}L_hH_h^{-1/2}).\label{rig:beta-matrix}
\end{equation}
Таким образом, одинаковая двойственная норма невязки может соответствовать
разным ошибкам при разных частотах. Для равномерного вывода нужна общая
положительная нижняя граница $\beta_h$ на рассматриваемом семействе.

\subsection{Возмущения сборки}
Пусть исходная система имеет вид $L_hU=F$, а вычисленная система ---
$(L_h+\Delta L_h)\widetilde U=F+\Delta F$. Обозначим
$\eta=\|H_h^{-1/2}\Delta L_hH_h^{-1/2}\|_2$, где $\|\cdot\|_2$ для
матрицы есть спектральная норма. При $\eta<\beta_h$ минимальное
сингулярное число нормированной возмущённой матрицы не меньше
$\beta_h-\eta$. Это следует из неравенства треугольника для
$(L_h+\Delta L_h)z$ в соответствующих нормах. Вычитая системы, получаем
правую часть $\Delta F-\Delta L_hU$ для разности решений и оценку
\begin{equation}
 \|\widetilde U-U\|_{H_h}
 \le\frac{\|\Delta F\|_{H_h^{-1}}+\eta\|U\|_{H_h}}{\beta_h-\eta}.
 \label{rig:operator-perturbation}
\end{equation}
Например, ошибка квадратуры одной только нагрузки входит через $\Delta F$.

\section{Переход от сетки к непрерывной задаче}
\label{rig:fem-convergence}

\begin{proposition}[Квазиоптимальность метода Галёркина]
Пусть $u_s\in V$ решает непрерывную задачу и $\beta_h>0$. Тогда при точном
интегрировании
\begin{equation}
 \|u_s-u_{s,h}\|_{V,k}\le
 \left(1+\frac{C_a(k)}{\beta_h(k)}\right)
 \inf_{v_h\in V_h}\|u_s-v_h\|_{V,k}.\label{rig:quasi-galerkin}
\end{equation}
\end{proposition}
\begin{proof}
Для любого $v_h\in V_h$ и любого теста $w_h$ имеем
$a_k(u_{s,h}-v_h,w_h)=a_k(u_s-v_h,w_h)$. Условия \eqref{rig:beta}
и \eqref{rig:continuity} дают
$\|u_{s,h}-v_h\|_{V,k}\le (C_a/\beta_h)\|u_s-v_h\|_{V,k}$.
Остаётся прибавить $\|u_s-v_h\|_{V,k}$ и взять инфимум.
\end{proof}
При фиксированной частоте для сходимости достаточно равномерной
устойчивости и условия $\inf_{v_h\in V_h}\|u_s-v_h\|_{V,k}\to0$ при $h\to0$.
В частности, на регулярных по форме треугольных сетках при $u_s\in H^2$
и максимальном диаметре элемента $h$ узловой интерполянт удовлетворяет
\begin{equation}
 \|u_s-I_hu_s\|_{V,k}\le C_I(h+kh^2)|u_s|_{H^2}.\label{rig:interpolation}
\end{equation}
Скачки свойств среды могут нарушить глобальную $H^2$-регулярность.
При изменении $k$ нужно учитывать также зависимость решения и $\beta_h$
от частоты.

Если фактические форма $a_h$ и нагрузка $\ell_h$ включают квадратуру,
возьмём любое приближение $v_h\in V_h$ и положим
$\psi_h(w_h)=\ell_h(w_h)-a_h(v_h,w_h)$. Пусть $\widetilde u_h$
решает фактическую задачу, а $\widetilde\beta_h>0$ есть её константа
устойчивости. Применяя \eqref{rig:residual-estimate} к
$\widetilde u_h-v_h$ и прибавляя ошибку аппроксимации, получаем
\begin{equation}
 \|u_s-\widetilde u_h\|_{V,k}\le\|u_s-v_h\|_{V,k}
                   +\widetilde\beta_h^{-1}\|\psi_h\|_{V_h',k}.
 \label{rig:consistency-residual}
\end{equation}
Здесь $v_h$ можно выбирать и при такой регулярности $u_s$, при которой
узловые значения самой функции не определены.

\section{Согласованная и диагональная матрицы массы}
\label{rig:mass-equivalence}

Для эрмитовых матриц запись $X\preceq Y$ означает
$z^*Xz\le z^*Yz$ при всех $z$. Связь между двумя метриками
главы~\ref{math:fem} особенно проста для линейных треугольников.

\begin{theorem}[Сравнение двух норм]
Для полной согласованной матрицы массы $M$ и $W=\diag(M\bm1)$ на
невырожденной треугольной сетке
\begin{equation}
 M\preceq W\preceq4M,\qquad
 \|z\|_M\le\|z\|_W\le2\|z\|_M.\label{rig:mass-bounds}
\end{equation}
\end{theorem}
\begin{proof}
На треугольнике площади $|T|$ имеем
$M_T=|T|(I+\bm1\bm1^\top)/12$, $W_T=|T|I/3$.
У постоянного вектора собственное значение $M_T$ равно $|T|/3$,
а у векторов с нулевой суммой координат оно равно $|T|/12$.
Следовательно, локально $M_T\preceq W_T\preceq4M_T$.
Суммирование локальных квадратичных форм по треугольникам даёт требуемые
неравенства, в том числе для комплексных $z$.
\end{proof}
Коэффициент 4 достигается уже на одном элементе. Результат относится
именно к полной линейной треугольной массе. Если имеются фиксированные
граничные значения, обе уже собранные формы можно ограничить на свободные
степени свободы. Повторное суммирование строк ограниченной $M$ даст другую
диагональную матрицу. Для элементов более высокого порядка нужны свои оценки.

\section{Наилучшее приближение в пространстве мод}
\label{rig:projection}

\begin{theorem}[Ортогональная проекция и ошибка коэффициентов]
Пусть $M=M^*>0$, $\mu\in\C^N$ и $\Phi\in\C^{N\times r}$ имеет
линейно независимые столбцы. Обозначим через $G=\Phi^*M\Phi$ матрицу
Грама и для $U\in\C^N$ определим
\begin{equation}
 Gc_*=\Phi^*M(U-\mu),\qquad P_MU=\mu+\Phi c_*,\qquad
 \rho=U-P_MU.\label{rig:projection-def}
\end{equation}
Тогда $P_MU$ есть единственное ближайшее к $U$ поле из
$\mu+\operatorname{ran}\Phi$, а при любом $c\in\C^r$
\begin{equation}
 \|U-\mu-\Phi c\|_M^2=\|\rho\|_M^2+\|c-c_*\|_G^2.
 \label{rig:projection-identity}
\end{equation}
\end{theorem}
\begin{proof}
Из независимости столбцов следует $x^*Gx=\|\Phi x\|_M^2>0$ при $x\ne0$.
Поэтому коэффициенты $c_*$ определены однозначно. По
\eqref{rig:projection-def} выполнено $\Phi^*M\rho=0$.
В разложении $U-\mu-\Phi c=\rho+\Phi(c_*-c)$ слагаемые
$M$-ортогональны, что даёт \eqref{rig:projection-identity}.
Его правая часть минимальна ровно при $c=c_*$.
\end{proof}
При $r=0$ пространство состоит из одного поля $\mu$ и остаётся только
остаток. Для раздельных вещественной и мнимой частей та же теорема
действует в $\R^{2N}$ с блочными матрицами $\diag(M,M)$ и
$\diag(\Phi_R,\Phi_I)$. При этом используют только независимые
активные столбцы каждой компоненты.

Пусть $P_W$ есть проекция на то же аффинное пространство в метрике $W$.
По минимальности обеих проекций и по \eqref{rig:mass-bounds}
\begin{equation}
 \|U-P_MU\|_M\le\|U-P_WU\|_M\le\|U-P_WU\|_W
 \le\|U-P_MU\|_W\le2\|U-P_MU\|_M.\label{rig:lumped-quasi}
\end{equation}
Таким образом, измеренная в $M$ ошибка $W$-реконструкции может превышать
оптимум. Её отношение к ненулевому оптимуму для этих элементов не больше двух.

\begin{worked}{Задача. Найдите обе проекции на одном треугольнике}
Пусть $|T|=1/2$, $U=(0,1,0)^\top$, а допустимые поля имеют вид
$(c,0,0)^\top$. Найдите коэффициент, дающий наименьшую ошибку в каждой
из норм.

\textbf{Решение.} Так как $W=I/6$, имеем $c_W=0$. Для полной матрицы
массы из \eqref{math:local-mass} уравнение проекции имеет вид
$(2/24)c_M=1/24$, откуда $c_M=1/2$. Квадрат $M$-ошибки для произвольного
$c\in\R$ равен $(2c^2-2c+2)/24$. При $c_W$ он равен $1/12$, при $c_M$
равен $1/16$. Соседние узловые функции перекрываются внутри элемента,
поэтому ненулевой коэффициент при первой функции улучшает приближение
второй в интегральной норме.
\end{worked}

\subsection{Почему SVD выбирает оптимальное пространство}
\label{rig:pod-optimality}
Для центрированных столбцовых снимков $S\in\C^{N\times m}$ и
$M=R^*R>0$ положим $Y=RS$. Любому $M$-ортонормированному базису
$\Psi\in\C^{N\times r}$ соответствует $Q=R\Psi$, причём $Q^*Q=I$.
Сумма квадратов ошибок проекции снимков равна
\[
 \|(I-QQ^*)Y\|_F^2=\operatorname{tr}(YY^*)-
                              \operatorname{tr}(Q^*YY^*Q).
\]
В собственном базисе $YY^*$ второй след равен $\sum_j\sigma_j^2p_j$.
Здесь спектр при необходимости дополнен нулями, а $p_j$ есть диагональные
элементы ортогонального проектора $QQ^*$ в этом базисе, так что
$0\le p_j\le1$ и $\sum_jp_j=r$. Максимум суммы достигается, когда весь
вес отведён первым $r$ собственным векторам. Следовательно, минимум
остатка равен $\sum_{j>r}\sigma_j^2$. Это доказывает оптимальность,
использованную в главе~\ref{math:pod}, и для комплексных снимков.
Если $m\ll N$, ненулевые квадраты сингулярных чисел находят по меньшей
матрице $Y^*Y=S^*MS$. Для нормированного собственного вектора $v_j$ и
$\sigma_j>0$ соответствующая мода равна $Sv_j/\sigma_j$.
Деление на очень малое $\sigma_j$ чувствительно к округлению.

\section{Чувствительность восстановления и сети}
\label{rig:decoder-stability}

Для декодера $D(c)=\mu+\Phi c$ квадрат нормы изменения поля равен
$\delta c^*G\delta c$. Поэтому при $r>0$
\begin{equation}
 \sqrt{\lambda_{\min}(G)}\,\|\delta c\|_2
 \le\|D(c+\delta c)-D(c)\|_M
 \le\sqrt{\lambda_{\max}(G)}\,\|\delta c\|_2.\label{rig:decoder-conditioning}
\end{equation}
При нормированных координатах $q$ и физическом векторе коэффициентов
$c=m_c+S_c q$ имеем
\begin{equation}
 \|D(m_c+S_c q_1)-D(m_c+S_c q_2)\|_M
 \le\underbrace{\|G^{1/2}S_c\|_2}_{\kappa_c}\|q_1-q_2\|_2.
 \label{rig:scaled-decoder}
\end{equation}
При $M$-ортонормированных модах $G=I$. Если ошибка затронула также
среднее и моды, $\widetilde D(c)=\mu+\delta\mu+(\Phi+\Delta\Phi)c$,
то из явной разности следует
\begin{equation}
 \|\widetilde D(c+\delta c)-D(c)\|_M
 \le\|\delta\mu\|_M+\|\Phi\delta c\|_M
 +\|\Delta\Phi\|_{2\to M}\|c+\delta c\|_2.
 \label{rig:decoder-perturb}
\end{equation}
Здесь $\|T\|_{2\to M}=\sup_{x\ne0}\|Tx\|_M/\|x\|_2$. Данная
чувствительность определяется пространством мод и масштабом коэффициентов.
Чувствительность решения краевой задачи задаёт другая константа, $\beta_h$.

\subsection{Как изменение признаков передаётся через сеть}
\label{rig:network-stability}
Для скрытых таблиц $Z$ используем норму Фробениуса. Один остаточный блок
действует как
\[
 F_\ell(Z)=Z+\alpha_\ell\sigma\bigl(ZW_{\rm loc}+\bm1 b^\top+
                                             B_\ell ZW_{\rm ch}\bigr).
\]
Пусть покомпонентная активация липшицева с константой $L_\sigma$.
Из $\|BZW\|_F\le\|B\|_2\|Z\|_F\|W\|_2$ и неравенства треугольника
следует
\begin{equation}
 \operatorname{Lip}(F_\ell)\le K_\ell=
 1+|\alpha_\ell|L_\sigma
 (\|W_{\rm loc}\|_2+\|B_\ell\|_2\|W_{\rm ch}\|_2).
 \label{rig:block-lip}
\end{equation}
Для синуса $L_\sigma=1$. Для точной GELU, $\sigma(t)=t\Phi_{\rm N}(t)$,
можно взять $L_\sigma=1+1/\sqrt{2\pi e}$, поскольку
$\sigma'(t)=\Phi_{\rm N}(t)+t\varphi_{\rm N}(t)$, где $\Phi_{\rm N}$ и
$\varphi_{\rm N}$ есть функция распределения и плотность стандартной
нормальной случайной величины.

В базовой архитектуре три общих признака повторяются в каждой из
$R_t=\max(r_{\rm out},r_A,r_n)$ строк. Каждый коэффициент материала
появляется один раз. Переход от плоского набора признаков к таблице
имеет норму $\sqrt{R_t}$. Если $W_{\rm in},W_{\rm out}$ есть матрицы
начального и выходного линейных слоёв, для сети из $L$ блоков имеем
\begin{equation}
 \operatorname{Lip}(N_\vartheta)\le\|W_{\rm out}\|_2
 \left(\prod_{\ell=1}^L K_\ell\right)
 L_\sigma\|W_{\rm in}\|_2\sqrt{R_t}.\label{rig:network-lip}
\end{equation}
Оценка учитывает активацию после входного слоя. Выбор первых выходных
строк имеет норму не больше единицы. Правая часть характеризует изменение
\emph{прогноза} между двумя входами. Для оценки его \emph{точности}
требуется сравнение с оператором решения.

\section{Аппроксимация оператора на семействе параметров}
\label{rig:identifiability}
\label{rig:total-error}

Определим модель на наборе параметров $\mathcal P$ формулой
\begin{equation}
 \widehat U(p)=D_h\!\left(m_c+S_cN_\vartheta(E_h(p))\right),\qquad
 p\in\mathcal P.\label{rig:factorized-model}
\end{equation}
Кодирование $E_h$ включает доступные сеточные признаки, их проекции
и нормировку. Детерминированная сеть выдаёт одинаковый ответ на два
запроса с одинаковым кодом. Если $E_h(p_1)=E_h(p_2)$, то для эталонов
$U_1,U_2$ на одной сетке
\begin{equation}
 \max_{j=1,2}\|U_j-\widehat U(p_j)\|_M\ge
                   \tfrac12\|U_1-U_2\|_M.\label{rig:input-floor}
\end{equation}
Действительно, расстояние между $U_1$ и $U_2$ не превосходит суммы
расстояний от них до их общего прогноза. Например, при нулевых входных
рангах два разных материала с одинаковыми $k,\theta$ получили бы один код.

\subsection{Разделение источников погрешности}
Пусть $c_*(p)$ и $\rho(p)$ есть коэффициенты и остаток точной $M$-проекции
эталона из \eqref{rig:projection-def}. Обозначим через $z(p)$ выбранные
идеальные признаки, через $\widehat z(p)$ --- фактически использованные
признаки. Для заданной функции $g$, сопоставляющей признакам нормированные
коэффициенты, положим
\begin{equation}
 a(p)=\|c_*(p)-m_c-S_c g(z(p))\|_G.\label{rig:representation-defect}
\end{equation}
Этот член учитывает точность самого описания через $g$ и $z$. При
неразличимых входах точное соответствие может оказаться невозможным.

\begin{proposition}[Оценка при заданных погрешностях звеньев]
Пусть $g$ липшицева с константой $L_g$ на множестве, содержащем
$z(p)$ и $\widehat z(p)$, и пусть в точке $p$ выполнены
$\|\widehat z-z\|_2\le\delta_z$,
$\|N_\vartheta(\widehat z)-g(\widehat z)\|_2\le\delta_N$.
Если вычисленный выход в нормированных координатах равен
$N_\vartheta(\widehat z(p))+e_q$, то для фиксированных среднего и мод
\begin{align}
 \|U(p)-\widehat U(p)\|_M
 &\le\sqrt{\|\rho(p)\|_M^2+Q(p)^2},\label{rig:end-to-end}\\
 Q(p)&=a(p)+\kappa_c\bigl(L_g\delta_z+\delta_N+\|e_q\|_2\bigr).\notag
\end{align}
\end{proposition}
\begin{proof}
В разности $\widehat c-c_*$ добавим и вычтем значения
$m_c+S_c g(z)$ и $m_c+S_c g(\widehat z)$. Неравенство треугольника,
условие Липшица и \eqref{rig:scaled-decoder} дают
$\|\widehat c-c_*\|_G\le Q(p)$. Остаётся применить точное тождество
\eqref{rig:projection-identity}.
\end{proof}
Если изменились также среднее и моды, к правой части добавляют
$\|\delta\mu\|_M+\|\Delta\Phi\|_{2\to M}\|\widehat c\|_2$. При
диагональной нормировке входов изменение одной координаты на $\delta y_j$
изменяет нормированную координату на $\delta y_j/s_j$. Поэтому малые
масштабы $s_j$ могут увеличить $\delta_z$.

Пусть $\widetilde u_h(p)$ есть конечно-элементная функция с эталонным
вектором $U(p)$, полученным при фактической сборке. Для перехода к
непрерывной задаче прибавляем $\|u_s(p)-\widetilde u_h(p)\|_{L^2}$
к оценке \eqref{rig:end-to-end}. Этот член включает погрешность
квадратуры. Его можно оценить по \eqref{rig:consistency-residual},
поскольку $\|v\|_{L^2}\le k^{-1}\|v\|_{V,k}$.
При сравнении с решением внешней задачи нужно учесть ещё ошибку
искусственной границы. Полученная величина относится к одному $p$.

\begin{idea}{Когда из оценки следует сходимость операторов}
Пусть все поля рассматриваются на общей области, а множество допустимых
параметров $\mathcal P$ фиксировано. Если при уточнении модели каждое
слагаемое приведённой абсолютной оценки стремится к нулю \emph{равномерно
по $p\in\mathcal P$}, то
$\sup_{p\in\mathcal P}\|u_s(p)-\widehat u_h(p)\|_{L^2}\to0$ для
усечённой задачи. Для равномерной относительной оценки достаточно также
границы $\inf_{p\in\mathcal P}\|u_s(p)\|_{L^2}>0$. Средняя ошибка на
конечном тестовом наборе характеризует только наблюдавшиеся запросы.
\end{idea}

\section{От математической нормы к вычисленной ошибке}
\label{rig:floating}
Вычисление \eqref{rig:projection-identity} в пространстве коэффициентов
позволяет заранее найти остатки и малую матрицу Грама. Для согласования
с фактическим декодером среднее и моды берут после округления до используемой
при восстановлении точности. Базовая реализация хранит их в одинарной
точности. Матрицу Грама, проекции и остатки по этим данным вычисляют
в двойной точности, причём остаток интегрируют непосредственно.
Это позволяет избежать вычитания двух почти равных энергий. При сравнении
с прямым восстановлением поля учитывают округления обеих процедур.

Если $e=\|U-\widehat U\|_M^2$ и $b=\|U\|_M^2$, базовая функция потерь
усредняет по примерам величину
\begin{equation}
 \ell_\varepsilon(e,b)=
 \frac{e}{\sqrt{\max(e,\varepsilon)}\sqrt{\max(b,\varepsilon)}},
 \qquad \varepsilon=10^{-24}.\label{rig:regularized-loss}
\end{equation}
При $e,b\ge\varepsilon$ это обычная относительная ошибка.
Вблизи $e=0$ такая запись даёт конечную производную по $e$ и сохраняет
нулевую потерю при точном прогнозе. Для почти нулевых полей результат
зависит от принятого правила. При чтении численных результатов нужно
уточнять, приведена ли ошибка по \eqref{math:relative-error} или значение
регуляризованной функции потерь.

\begin{worked}{Задача. Два запроса с одинаковым кодом}
Пусть два входа получили одинаковый код, а их эталонные поля равны
$\phi$ и $3\phi$, где $\|\phi\|_M=1$. Общий ответ имеет вид $t\phi$,
$t\in\R$. Какой ответ минимизирует наибольшую абсолютную ошибку?
Как изменится ответ, если ошибки считать относительно каждого эталона?

\textbf{Решение.} При $1\le t\le3$ абсолютные ошибки равны $t-1$ и
$3-t$. За пределами этого отрезка обе ошибки уменьшаются при движении
к нему. На отрезке максимум минимален при равенстве ошибок, то есть при
$t=2$. Он равен единице, что совпадает с нижней границей
\eqref{rig:input-floor}. Относительные ошибки равны $t-1$ и $(3-t)/3$.
Их равенство даёт $t=3/2$ и максимум $1/2$. Таким образом, переход к
относительной ошибке меняет оптимальный общий прогноз даже при тех же
полях и том же пространстве мод.
\end{worked}
